
% ============================================================
% DOKUMENTKLASSE
% ============================================================

% \documentclass legt die Grundart des Dokuments fest.
% "article" ist für kürzere wissenschaftliche Arbeiten,
% Berichte und Dokumentationen geeignet.
\documentclass{article}


% ============================================================
% PAKETE
% ============================================================

% inputenc ermöglicht die Verwendung von UTF-8-Zeichen,
% z. B. ä, ö, ü und ß.
% Bei aktuellen LaTeX-Versionen ist UTF-8 bereits Standard,
% der Befehl kann aber zur Verdeutlichung stehen bleiben.
\usepackage[utf8]{inputenc}

% fontenc sorgt dafür, dass Sonderzeichen korrekt
% in der verwendeten Schrift kodiert werden.
% T1 ist für europäische Sprachen geeignet.
\usepackage[T1]{fontenc}

% babel passt das Dokument an die deutsche Sprache an.
% Dadurch werden z. B. Überschriften und Trennungen
% automatisch auf Deutsch angepasst.
\usepackage[ngerman]{babel}

% amsmath stellt zusätzliche Funktionen für mathematische
% Formeln und mathematische Umgebungen zur Verfügung.
\usepackage{amsmath}

% amssymb stellt zusätzliche mathematische Symbole bereit.
\usepackage{amssymb}

% geometry ermöglicht die Einstellung der Seitenränder.
% a4paper = Papierformat DIN A4
% margin=2.5cm = 2,5 cm Rand auf allen Seiten
\usepackage[a4paper,margin=2.5cm]{geometry}


% ============================================================
% BEGINN DES DOKUMENTS
% ============================================================

% \begin{document} markiert den Beginn des eigentlichen
% Dokuments. Alles davor gehört zum sogenannten Vorspann.
\begin{document}


% ============================================================
% TITELBEREICH
% ============================================================

% \title definiert den Titel des Dokuments.
\title{Analyse von Datenbank-Performance, Streuung und Indizierung}

% \author definiert den Autor.
\author{David Rinner}

% \date definiert das Datum.
% \today setzt automatisch das aktuelle Datum ein.
\date{\today}

% \maketitle erzeugt den zuvor definierten Titel,
% Autor und das Datum im Dokument.
\maketitle


% ============================================================
% ABSCHNITT 1
% ============================================================

% \section erstellt eine nummerierte Überschrift.
\section{Einleitung und Methodik}

Ziel dieser Analyse ist die Untersuchung des Verhaltens von

SQLite-Datenbanken bei großen Datenmengen
% $ ... $ öffnet und schließt den sogenannten Mathematikmodus.
% \geq bedeutet "größer oder gleich".
% \, erzeugt einen kleinen horizontalen Abstand.
% 500\,000 sorgt dafür, dass 500 und 000 nicht unschön
% auseinandergerissen werden.
($\geq 500\,000$ Datensätzen).

Dabei wurden die Streuung von Zufallsdaten, die mithilfe von Python und

% \texttt{...} stellt Text in einer Schreibmaschinenschrift dar.
% Besonders geeignet für Programmcode, Dateinamen und Befehle.
\texttt{Faker} generiert wurden, die Auswirkungen von Indizes auf die

Performance sowie der Einfluss eines Daten-Bias untersucht.


% ============================================================
% ABSCHNITT 2
% ============================================================

\section{Streuungsanalyse der Zufallsdaten}

Da Namen keine Distanz im herkömmlichen Sinne besitzen, lässt sich ihre

Streuung nicht direkt über geometrische Abstände bestimmen.

Stattdessen wird die Streuung über die
% \textbf{...} stellt den enthaltenen Text fett dar.
\textbf{Häufigkeitsverteilung der Vornamen} definiert.

Jeder Vorname wird dabei als Kategorie betrachtet, deren Häufigkeit über

alle Datensätze hinweg gezählt wird.


% ============================================================
% AUFZÄHLUNG
% ============================================================

% \begin{itemize} beginnt eine Aufzählung.
\begin{itemize}

    % \item erzeugt einen einzelnen Aufzählungspunkt.

    \item \textbf{Durchschnittliche Häufigkeit:}
    $\approx 249{,}75$

    Vorkommen pro Vorname.

    % \approx bedeutet "ungefähr gleich".
    % {,} sorgt dafür, dass das Komma im Mathematikmodus
    % als normales Komma dargestellt wird.

    \item \textbf{Varianz der Häufigkeiten:}
    $267{,}28$

    (Standardabweichung $\sigma \approx 16{,}35$).

    % \sigma erzeugt den griechischen Buchstaben σ.
    % \approx erzeugt das Zeichen ≈.

% \end{itemize} beendet die Aufzählung.
\end{itemize}

Die geringe Varianz zeigt, dass die verwendete Bibliothek

\texttt{Faker} die Vornamen relativ gleichmäßig über die Datensätze

verteilt.


% ============================================================
% ABSCHNITT 3
% ============================================================

\section{Performance und Speicherbedarf mit und ohne Index}

Ohne Index muss die Datenbank bei einer Abfrage

die gesamte Tabelle durchsuchen und benötigt in der vorliegenden

Messung etwa $0{,}178\,\mathrm{s}$.

% \mathrm{...} sorgt dafür, dass "s" als normale Einheit
% und nicht als mathematische Variable dargestellt wird.
%
% $0{,}178\,\mathrm{s}$
% bedeutet:
% 0,178 Sekunden.


Nach der Erstellung eines Indexes

sinkt die Abfragezeit auf etwa $0{,}001\,\mathrm{s}$.
Dadurch wird die Abfrage in diesem Beispiel deutlich beschleunigt.


% ============================================================
% SPEICHERBEDARF
% ============================================================

Dieser Geschwindigkeitsgewinn geht jedoch mit einem erhöhten

Speicherbedarf einher. Während die Datenbank vor der Erstellung des

Indexes

$11\,755\,520\,\mathrm{Byte}$

beanspruchte, steigt der Speicherbedarf nach dem Aufbau des

B-Tree-Indexes auf

$19\,578\,880\,\mathrm{Byte}$.

% \, erzeugt einen kleinen Abstand zwischen Zahlenblöcken.
%
% Beispiel:
% 11\,755\,520
%
% wird typografisch besser dargestellt als:
% 11 755 520


Der zusätzliche Speicherbedarf entsteht dadurch, dass neben der

eigentlichen Tabelle eine zusätzliche Indexstruktur gespeichert werden

muss.


% ============================================================
% ABSCHNITT 4
% ============================================================

\section{Untersuchung der Biased Database}

Wird die Verteilung der Vornamen gezielt verändert, sodass

50\,\% aller Datensätze denselben Vornamen (\texttt{Simon}) tragen,

verändert sich die Geschwindigkeit der Suche.

% \% erzeugt das Prozentzeichen.
% Das einfache "%" wäre in LaTeX ein Kommentarzeichen.
%
% Deshalb muss man schreiben:
% 50\,\%
%
% \, = kleiner Abstand
% \% = Prozentzeichen


% ============================================================
% AUFZÄHLUNG
% ============================================================

\begin{itemize}

    \item \textbf{Seltener Wert (\texttt{Peter}):}

    Wird nach einem selten vorkommenden Namen gesucht, kann der Index

    die passenden Datensätze sehr schnell finden. In der Messung wurde

    eine Abfragezeit von etwa $0{,}001\,\mathrm{s}$ erreicht.


    \item \textbf{Häufiger Wert (\texttt{Simon}):}

    Da \texttt{Simon} in 50\,\% der Datensätze vorkommt, muss die

    Datenbank sehr viele passende Datensätze finden. Dadurch ist die

    Suche mit dem Index langsamer als bei einem seltenen Namen.

    Mit Index benötigt die Abfrage etwa $0{,}050\,\mathrm{s}$,

    während ohne Index etwa $0{,}131\,\mathrm{s}$ benötigt werden.

\end{itemize}


% ============================================================
% ERKLÄRUNG DES VERHALTENS
% ============================================================

% \textbf{} wird verwendet, um eine Überschrift innerhalb
% eines normalen Absatzes fett darzustellen.
\textbf{Grund für dieses Verhalten:}

Der Unterschied liegt darin, wie oft ein bestimmter Name in der

Datenbank vorkommt. Bei einem seltenen Namen wie \texttt{Peter} gibt

es nur wenige passende Datensätze. Der Index kann diese daher schnell

finden.

Bei \texttt{Simon} sieht es anders aus: Der Name kommt in der Hälfte

aller Datensätze vor. Bei insgesamt $500\,000$ Datensätzen sind das

ungefähr $250\,000$ Treffer. Der Index muss deshalb deutlich mehr

passende Einträge verarbeiten.

Trotzdem ist die Suche mit dem Index in diesem Beispiel noch schneller

als die Suche ohne Index. Das zeigt, dass ein Index besonders bei

selten vorkommenden Werten hilfreich ist. Bei sehr häufig vorkommenden

Werten fällt der Vorteil dagegen kleiner aus.


% ============================================================
% ABSCHNITT 5
% ============================================================

\section{Fazit}

Die Untersuchung zeigt, dass die Verwendung eines Indexes bei großen

Datenmengen einen erheblichen Einfluss auf die Abfragegeschwindigkeit

haben kann. Besonders bei Suchbegriffen

mit selten vorkommenden Werten ist der Geschwindigkeitsgewinn deutlich.

Gleichzeitig erhöht ein Index den Speicherbedarf der Datenbank, da

zusätzliche Strukturen angelegt und verwaltet werden müssen.

Die Untersuchung der \texttt{bias.db} zeigt außerdem, dass die

Verteilung der Daten einen wesentlichen Einfluss auf die Effektivität

eines Indexes hat. Bei häufig vorkommenden Werten ist der Vorteil eines

Indexes geringer als bei seltenen Werten.

Damit wird deutlich, dass bei der Optimierung von Datenbanken nicht

nur die Anzahl der Datensätze, sondern auch deren Verteilung und die

Selektivität der verwendeten Suchbedingungen berücksichtigt werden

müssen.


% ============================================================
% ENDE DES DOKUMENTS
% ============================================================

% \end{document} beendet das LaTeX-Dokument.
% Alles danach wird von LaTeX nicht mehr zum Dokument gezählt.
\end{document}

